\subsection{Dual vector spaces}

Let \(\mathbf{F}\) and \(\mathbf{G}\) be two real vector spaces and let \((x, y) \mapsto \mathbf{B}(x, y)\) be a bilinear form on \(\mathbf{F} \times \mathbf{G}\) . We say that the bilinear form \(\mathbf{B}\) puts the vector spaces \(\mathbf{F}\) and \(\mathbf{G}\) in duality, or that F and G are in duality (relative to B). Recall that we say that \(x \in F\) and \(y \in G\) are orthogonal (for the duality defined by B) if \(\mathbf{B}(x, y) = 0\) ; we say that a subset M of F and a subset N of G are orthogonal if every \(x \in M\) is orthogonal to every \(y \in N\) (A, IX, § 1.2).

We say that the duality defined by B is separating in F (resp. in G) if it satisfies the following condition :

(D₁) For every \(x \neq 0\) in F, there exists \(y \in G\) such that \(\mathrm{B}(x, y) \neq 0\) . (resp. (D \(_{II}\) ) For every y ≠ 0 in G, there exists x ∈ F such that B(x, y) ≠ 0.)

The duality defined by B is said to be separating if it is both separating in F and in G. For this to be so, it is necessary and sufficient that the bilinear form B should be separating in the sense of A, IX, § 1.1. More precisely we have the following result:

PROPOSITION 1. --- Let F, G be two real vector spaces and B a bilinear form on F × G. Let

\[
d _ {\mathrm{B}}: y \mapsto \mathrm{B} (., y),
\]

\[
s _ {\mathbf {B}} \colon x \mapsto \mathrm{B} (x,.)
\]

be linear mappings of G in the dual F* of F and of F in the dual G* of G, associated respectively to the right and to the left of B (A, IX, § 1.1). Then B puts F and G in a duality separating in G (resp. in F), if and only if \(d_{\mathbf{B}}\) (resp. \(s_{\mathbf{B}}\) ) is injective.

When F and G are put in separating duality by B, we often identify F (resp. G) with a subspace of \(G^{*}\) (resp. \(F^{*}\) ) by means of \(s_{B}\) (resp. \(d_{B}\) ). When we consider F (resp. G) as a subspace of \(G^{*}\) (resp. \(F^{*}\) ) without specifying how this identification is to be made, we are always using the preceding identifications; the bilinear form B is then identified with the restriction to \(F \times G\) of the canonical bilinear form:

\[
(x ^ {*}, x) \mapsto \langle x, x ^ {*} \rangle \quad (\text { resp. } (x, x ^ {*}) \mapsto \langle x, x ^ {*} \rangle).
\]

Examples. --- 1) Let E be a vector space and let E* be its dual. The canonical bilinear form \((x, x^{*}) \mapsto \langle x, x^{*} \rangle\) on \(E \times E^{*}\) (A, II, §2.3) puts E and \(E^{*}\) in separating duality: for \((\mathrm{D}_{\mathrm{II}})\) is true because of the definition of the relation \(x^{*} \neq 0\) , and we know on the other hand, that for all \(x \neq 0\) in E, there exists a linear form \(x^{*} \in E^{*}\) such that \(\langle x, x^{*} \rangle \neq 0\) (A, II, §7.5, th. 6), which proves \((\mathrm{D}_{\mathrm{I}})\) ; the identifications of E with a subspace of \(E^{**}\) is made here by the canonical mapping \(c_{E}\) (loc. cit.).

When E is of finite dimension, the only subspace G of \(E^{*}\) that is in separating duality by the restriction to \(E \times G\) of the canonical bilinear form, is the space \(E^{*}\) itself; for, E being then canonically identified with \(E^{**}\) (loc. cit.), if we had \(G \neq E^{*}\) , there would exist \(a \neq 0\) in E such that \(\langle a, x^{*} \rangle = 0\) for all \(x^{*} \in G\) (A, II, §7.5, th. 7), which contradicts the hypothesis.

\begin{enumerate}
\def\labelenumi{\arabic{enumi})}
\setcounter{enumi}{1}
\tightlist
\item
  When E is an infinite dimensional vector space, and \(E'\) is a vector subspace of \(E^{*}\) , the duality between E and \(E'\) defined by the restriction to \(E \times E'\) of the canonical bilinear form is always separating in \(E'\) ; it can be separating in E even if \(E' \neq E^{*}\) . The most important case occurs where E is a topological vector space.
\end{enumerate}

DEFINITION 1. --- By the dual of a topological vector space E, we mean the subspace E' of E*, the dual of the vector space E, formed by the continuous linear forms on E.

When E is a Hausdorff locally convex space, the duality between E and its dual \(E'\) is separating: this follows from the Hahn-Banach theorem (II, p.~24, cor. 1) that for every \(x \neq 0\) in E, there exists \(x' \in E'\) such that \(\langle x, x' \rangle \neq 0\) .

Remarks. --- 1) When E is a topological vector space, the dual E* of the vector space E will be called the algebraic dual of E to avoid confusion. We note also that E* is the dual of the topological vector space obtained by giving E the finest locally convex topology (II, p.~25, Example 2).

\begin{enumerate}
\def\labelenumi{\arabic{enumi})}
\setcounter{enumi}{1}
\item
  The dual \(\mathbf{E}'\) of a topological vector space does not itself carry a topology, unless this is expressly stated.
\item
  If \(\mathbf{F}\) and \(\mathbf{G} \subset \mathbf{F}^*\) are in separating duality by the canonical bilinear form, then this is also true of \(\mathbf{F}\) and \(\mathbf{G}_1\) , for every subspace \(\mathbf{G}_1\) of \(\mathbf{F}^*\) such that \(\mathbf{G} \subset \mathbf{G}_1\) .
\end{enumerate}

